\documentclass[11pt,a4paper]{amsart}

\usepackage[margin=1in]{geometry}
\usepackage{amsmath,amssymb,amsthm,mathtools}
\usepackage[expansion=false]{microtype}
\usepackage[colorlinks=true,linkcolor=black,citecolor=black,urlcolor=black]{hyperref}
\usepackage[nameinlink,capitalize]{cleveref}
\hypersetup{
  pdftitle={Coefficientwise log-concavity of the cubic multiple-Pochhammer series},
  pdfauthor={John Shields},
  pdfkeywords={coefficientwise log-concavity, Turan inequality, Pochhammer symbol,
    log-concave sequence, beta-binomial distribution, Bernstein basis,
    Schur-concavity}
}
\pdfinfo{
  /ORCID (https://orcid.org/0009-0004-3882-4571)
  /DOI (10.5281/zenodo.21475821)
}

\numberwithin{equation}{section}

\theoremstyle{plain}
\newtheorem{theorem}{Theorem}[section]
\newtheorem{proposition}[theorem]{Proposition}
\newtheorem{lemma}[theorem]{Lemma}
\newtheorem{corollary}[theorem]{Corollary}

\theoremstyle{definition}
\newtheorem{definition}[theorem]{Definition}
\newtheorem{remark}[theorem]{Remark}

\newcommand{\E}{\mathbb E}
\newcommand{\Pp}{\mathbb P}
\newcommand{\Tur}{\Delta}

\title[Coefficientwise log-concavity of the cubic multiple-Pochhammer series]
{Coefficientwise log-concavity\\
of the cubic multiple-Pochhammer series}

\author{John Shields}
\date{July 28, 2026}
\thanks{Independent researcher, Ginza 2-15-2, KR Ginza II 5F, Chuo-ku, Tokyo,
Japan 104-0061.  Corresponding author:
\href{mailto:math@johnnyshields.com}{math@johnnyshields.com}.
ORCID: \href{https://orcid.org/0009-0004-3882-4571}{0009-0004-3882-4571}.
This work is archived at \href{https://doi.org/10.5281/zenodo.21475821}{doi:10.5281/zenodo.21475821}.}

\subjclass[2020]{Primary 26A51; Secondary 05A20, 33C20, 60E15}
\keywords{coefficientwise log-concavity, Tur\'an inequality, Pochhammer symbol,
log-concave sequence, beta-binomial distribution, Bernstein basis,
Schur-concavity}

\begin{document}

\begin{abstract}
Let \((f_n)_{n\ge1}\) be a nonnegative log-concave sequence with no internal
zeros, and define the cubic multiple-Pochhammer series
\(F_f(\mu;x)=\sum_{n\ge1}f_n\,(\mu)_{3n}\,x^n/(3n-1)!\).
Karp and Zhang established coefficientwise log-concavity of the analogous
series for multiplicity \(r=2\) and gave numerical counterexamples for \(r=4\), leaving the
cubic case open as their Conjecture~1.  We prove it: for all
\(\mu,\alpha,\beta\ge0\), every coefficient of the generalized Tur\'anian
\(F_f(\mu+\alpha;x)F_f(\mu+\beta;x)-F_f(\mu;x)F_f(\mu+\alpha+\beta;x)\) is
nonnegative.  Equivalently, the map \((f_n)\mapsto F_f\) carries every
such sequence to a coefficientwise log-concave family.  The proof is elementary
and noncomputational.  The degree-\(m\) coefficient of each product
\(F_f(u;x)F_f(v;x)\) is represented as a beta-binomial expectation, reducing
the theorem to the monotonicity of a weighted residue-class binomial kernel.
The constant-weight case rests on a Bernstein-basis positivity certificate
obtained from a root-of-unity filter; monotone weights are then inserted through
a one-sign-change decomposition of the complementary residue blocks.
A coefficientwise strictness criterion follows, as does ordinary log-concavity
in \(\mu\) on \((0,\infty)\) for \((f_n)\not\equiv0\) at every positive \(x\)
within the radius of convergence.
\end{abstract}

\maketitle

\section{Introduction}

A formal power series depending on a parameter \(\mu\ge0\) is called
\emph{coefficientwise log-concave} if, for every \(\mu,\alpha,\beta\ge0\), all
coefficients of the generalized Tur\'anian
\[
 F(\mu+\alpha;x)F(\mu+\beta;x)
 -F(\mu;x)F(\mu+\alpha+\beta;x)
\]
are nonnegative.  This property is stronger than the corresponding pointwise
inequality: whenever the relevant series converge at a positive value of \(x\), summing
the coefficient inequalities yields the generalized Tur\'an inequality for
\(\mu\mapsto F(\mu;x)\) there, simultaneously for all admissible shifts.  It has been
studied systematically for reciprocal-gamma, gamma-ratio,
hypergeometric, and Fox--Wright series; see
\cite{KalmykovKarp2013Reciprocal,KalmykovKarp2013Ratios,KarpSitnik2010,KarpZhang2024}.
Tur\'an-type inequalities for hypergeometric and confluent hypergeometric
functions have also been studied extensively
\cite{Baricz2008,BarnardGordyRichards2009}.

Karp and Zhang considered, for an integer \(r\ge2\), the modified
multiple-Pochhammer series
\begin{equation}\label{eq:general-r}
 F_{f,r}(\mu;x)
 =\sum_{n\ge1}f_n\frac{(\mu)_{rn}}{(rn-1)!}x^n,
\end{equation}
where \((f_n)\) is nonnegative and log-concave.  Karp and Zhang report that
A.\ Ninh suggested this construction to them in a private communication in
2017 \cite{KarpZhang2024}; related published work on log-concavity and
Tur\'an-type inequalities appears in \cite{NinhPham2025}.  They proved
coefficientwise log-concavity for \(r=2\) \cite[Thm.~2]{KarpZhang2024} and
gave numerical counterexamples for \(r=4\) \cite[Rem.~5]{KarpZhang2024},
leaving the cubic case \(r=3\) open as their Conjecture~1.  This note proves it.

Recall that a nonnegative sequence \((f_n)_{n\ge1}\) is log-concave with no
internal zeros if \(f_n^2\ge f_{n-1}f_{n+1}\) whenever the three terms are
defined, and its positive support is an interval of integers.  A nontrivial such
sequence is a P\'olya frequency sequence of order two, or \(\mathrm{PF}_2\), also
called doubly positive.  The interval-support condition does not follow from the
inequality, and \Cref{thm:main} fails without it (\Cref{rem:internal-zeros}).

Karp and Zhang build the inequality and the interval-support condition into
their definition of log-concavity \cite[Def.~1]{KarpZhang2024}.  They index the
ambient sequence from \(n=0\), although the modified series itself begins at
\(n=1\).  The log-concavity hypothesis is stated here directly for the
positive-index tail \((f_n)_{n\ge1}\).  The two formulations are equivalent:
restriction preserves log-concavity and interval support, while any admissible
tail extends to \(n=0\) by setting \(f_0=0\).  The hypotheses of \Cref{thm:main}
are therefore those of their Conjecture~1.

Except in \Cref{cor:ordinary}, we work with formal power series, so no growth
condition on \((f_n)\) is needed.

\begin{theorem}[Cubic multiple-Pochhammer theorem]\label{thm:main}
Let \((f_n)_{n\ge1}\) be a nonnegative log-concave sequence with no internal
zeros, and put
\begin{equation}\label{eq:F-def}
 F_f(\mu;x)=\sum_{n\ge1}f_n\frac{(\mu)_{3n}}{(3n-1)!}x^n.
\end{equation}
For every \(\mu,\alpha,\beta\ge0\), the formal power series
\begin{equation}\label{eq:Turan-def}
 \Tur_f(\mu;\alpha,\beta;x)
 =F_f(\mu+\alpha;x)F_f(\mu+\beta;x)
 -F_f(\mu;x)F_f(\mu+\alpha+\beta;x)
\end{equation}
has nonnegative coefficients.
\end{theorem}

\Cref{thm:main} concerns the transform \((f_n)\mapsto F_f\), not a single
function: every nonnegative log-concave sequence is sent to a parameter family
whose generalized Tur\'anians have nonnegative coefficients.  Summing at
\(\alpha=\beta\) and invoking continuity yields ordinary log-concavity of
\(\mu\mapsto F_f(\mu;x)\) where the series converges (\Cref{cor:ordinary}).  The
cubic case is the intervening multiplicity between the positive case \(r=2\) and
the first known failure at \(r=4\), and the residue-class positivity of
\Cref{sec:kernel} is the mechanism behind it.

For \(f_n\equiv1\), the triplication formula for the Pochhammer symbol gives
\[
 F_f(\mu;x)
 =3x\,\frac{\mathrm d}{\mathrm dx}\,
 {}_3F_2\!\left(\tfrac{\mu}{3},\tfrac{\mu+1}{3},\tfrac{\mu+2}{3};\,
 \tfrac13,\tfrac23;\,x\right),
\]
so the theorem yields coefficientwise generalized Tur\'an inequalities for the
differentiated \({}_3F_2\) family above; a general log-concave \((f_n)\) gives a
coefficient-weighted extension.

The proof has three components.  The degree-\(m\) coefficient of each product
\(F_f(u;x)F_f(v;x)\) is rewritten as a beta-binomial expectation
(\Cref{sec:beta-binomial}).  Two elementary
monotonicity lemmas (\Cref{sec:monotonicity-lemmas}) then supply the tools:
beta laws with a fixed sum of parameters are ordered by their imbalance against
symmetric functions increasing toward the midpoint, and a one-sign-change
weighting principle carries positivity through nondecreasing weights.  The
weighted residue-class binomial kernel is shown to have the required
monotonicity (\Cref{sec:kernel}); here complementary residue blocks have a
single sign change, and the unweighted sum carries a Bernstein-basis positivity
certificate.  This kernel theorem does not depend on the beta-binomial setup.  The
components are combined in \Cref{sec:proof}.

\section{Beta-binomial representation}\label{sec:beta-binomial}

This section rewrites the coefficient convolution underlying the Tur\'anian
\eqref{eq:Turan-def} as a beta-binomial expectation.

For \(m\ge2\), put
\begin{equation}\label{eq:w-from-f}
 w_k=f_kf_{m-k},\qquad 1\le k\le m-1.
\end{equation}
Log-concavity of \((f_n)\) implies
\[
 w_{k+1}=f_{k+1}f_{m-k-1}\ge f_kf_{m-k}=w_k
 \qquad\Bigl(1\le k<\bigl\lfloor\tfrac m2\bigr\rfloor\Bigr),
\]
the standard cross-product consequence \(f_{a+1}f_b\ge f_af_{b+1}\) for \(a<b\) of
log-concavity with no internal zeros.  With the symmetry \(w_k=w_{m-k}\), the
weights \eqref{eq:w-from-f} satisfy the hypotheses of \Cref{thm:kernel}.

For \(u,v\ge0\), the degree-\(m\) coefficient of the product \(F_f(u;x)F_f(v;x)\) is
\begin{equation}\label{eq:C-def}
 C_{m,f}(u,v)
 =\sum_{k=1}^{m-1}w_k
 \frac{(u)_{3k}(v)_{3(m-k)}}{(3k-1)![3(m-k)-1]!},
\end{equation}
with \(C_{m,f}(0,v)=0\) since \((0)_{3k}=0\).  The beta-binomial
representation below is used for \(u,v>0\).
Let \(N=3m\) and let \(J\) have the beta-binomial distribution
\[
 \Pp(J=j)=\binom Nj\frac{(u)_j(v)_{N-j}}{(u+v)_N},
 \qquad 0\le j\le N.
\]
For \(1\le j\le N-1\),
\[
 \frac{(u)_j(v)_{N-j}}{(j-1)!(N-j-1)!}
 =\frac{(u+v)_N}{N!}\,
 j(N-j)\Pp(J=j).
\]
Defining a weight on all integers \(0\le j\le N\) by
\[
 \widetilde w_j=
 \begin{cases}
  w_{j/3},&0<j<N\ \text{and}\ 3\mid j,\\
  0,&\text{otherwise},
 \end{cases}
\]
we obtain
\begin{equation}\label{eq:C-beta-binomial}
 C_{m,f}(u,v)
 =\frac{(u+v)_N}{N!}
 \E\!\left[J(N-J)\widetilde w_J\right].
\end{equation}
The factor \(J(N-J)\) vanishes at \(j=0\) and \(j=N\), so the endpoint terms do not
contribute.

Use the standard mixture representation \(P\sim\operatorname{Beta}(u,v)\) and
\(J\mid P\sim\operatorname{Bin}(N,P)\).  The identity
\[
 j(N-j)\binom Nj=N(N-1)\binom{N-2}{j-1}
\]
then transforms \eqref{eq:C-beta-binomial} into
\begin{equation}\label{eq:C-beta}
 C_{m,f}(u,v)
 =\frac{(u+v)_N}{N!}N(N-1)\E\,G_{m,w}(P),
\end{equation}
where the kernel
\begin{equation}\label{eq:G-weighted}
 G_{m,w}(p)
 =\sum_{k=1}^{m-1}w_k\binom{3m-2}{3k-1}
 p^{3k}(1-p)^{3(m-k)},
 \qquad 0\le p\le1,
\end{equation}
is defined for arbitrary symmetric nonnegative weights \(w_k=w_{m-k}\) and is
symmetric under \(p\mapsto1-p\).  The congruence \(j\equiv0\pmod3\) becomes
\(j-1\equiv2\pmod3\), which explains the binomial coefficient
\(\binom{3m-2}{3k-1}\).

By \eqref{eq:C-beta} and \Cref{lem:beta-order}, \Cref{thm:main} reduces to the
monotonicity of \(G_{m,w}\) on \([0,1/2]\) (\Cref{sec:kernel}).

\section{Two monotonicity lemmas}\label{sec:monotonicity-lemmas}

This section records two elementary monotonicity facts.  The weighting principle
for the paired residue blocks of \Cref{sec:kernel} is due to
\cite[Lem.~6]{KalmykovKarp2013Ratios}, proved in
\cite[Lem.~2.1]{KalmykovKarp2013Reciprocal}, and was applied to multiplicity two
in \cite{KarpZhang2024}.  The ordering of beta laws by imbalance used in
\Cref{sec:proof} is a fixed-sum likelihood-ratio comparison for folded beta
laws, obtained in \cite[Thm.~3]{BerenhautBergen2011} under the restriction that
each parameter gap be \(0\) or at least \(1\), and proved below without that
restriction.

\begin{lemma}\label{lem:weighting}
Let \(0\le w_1\le w_2\le\cdots\le w_L\) and suppose the real sequence
\(A_1,\dots,A_L\) has no more than one sign change, from nonpositive
to nonnegative.  If
\[
 \sum_{k=1}^L A_k\ge0,
\]
then
\[
 \sum_{k=1}^L w_kA_k\ge0.
\]
If in addition \(\sum_{k=1}^LA_k>0\), \(w_L>0\), and \(A_L>0\), then
\(\sum_{k=1}^Lw_kA_k>0\).
\end{lemma}

\begin{proof}
The claim is immediate if all \(A_k\) are nonnegative.  Otherwise, let \(k_0\) be the
largest index for which \(A_{k_0}<0\).  Put \(c=w_{k_0}\).  For \(k\le k_0\) we have
\(w_k\le c\) and \(A_k\le0\), while for \(k>k_0\) we have \(w_k\ge c\) and
\(A_k\ge0\).  Therefore
\[
 \sum_{k=1}^L w_kA_k
 \ge c\sum_{k=1}^L A_k\ge0.
\]
For the strict statement, assume \(\sum_kA_k>0\), \(w_L>0\), and \(A_L>0\).  If every
\(A_k\) is nonnegative, then \(\sum_kw_kA_k\ge w_LA_L>0\).  Otherwise \(A_L>0\) forces
\(k_0<L\).  If \(c>0\), then \(\sum_kw_kA_k\ge c\sum_kA_k>0\).  If \(c=0\), then
\(w_k=0\) for all \(k\le k_0\), so \(\sum_kw_kA_k=\sum_{k>k_0}w_kA_k\ge w_LA_L>0\).
\end{proof}

\begin{lemma}\label{lem:beta-order}
Fix \(s>0\).  For \(0\le d<s/2\), let
\(P_d\sim\operatorname{Beta}(s/2+d,\,s/2-d)\).
If \(H:[0,1]\to\mathbb R\) is symmetric, \(H(p)=H(1-p)\), and nondecreasing on
\([0,1/2]\), then
\begin{equation}\label{eq:beta-order}
 d_1\le d_2
 \quad\Longrightarrow\quad
 \E H(P_{d_1})\ge \E H(P_{d_2}).
\end{equation}
If \(H\) is strictly increasing on \((0,1/2)\) and \(d_1<d_2\), then the inequality
is strict.
\end{lemma}

\begin{proof}
The case \(d_1=d_2\) is immediate, so assume \(d_1<d_2\).  Put
\(Q_d=P_d(1-P_d)\in(0,1/4)\).  For
\[
 p_\pm(q)=\frac{1\pm\sqrt{1-4q}}{2},
 \qquad
 \ell(q)=\log\frac{p_+(q)}{p_-(q)},
\]
the density of \(Q_d\) is
\[
 g_d(q)=\frac{2}{B(s/2+d,s/2-d)}\,
 q^{s/2-1}(1-4q)^{-1/2}\cosh\!\bigl(d\ell(q)\bigr).
\]
The densities \(g_{d_1}\) and \(g_{d_2}\) share the factor
\(q^{s/2-1}(1-4q)^{-1/2}\), so
\[
 \frac{g_{d_2}(q)}{g_{d_1}(q)}
 =\frac{B(s/2+d_1,s/2-d_1)}{B(s/2+d_2,s/2-d_2)}\,
  \frac{\cosh(d_2\ell(q))}{\cosh(d_1\ell(q))}.
\]
The function \(\ell\) is strictly decreasing, since \(p_+(q)/p_-(q)\) decreases to
\(1\) as \(q\uparrow1/4\).  For \(d_2>d_1\), \(\cosh(d_2\ell)/\cosh(d_1\ell)\) is
strictly increasing in \(\ell>0\), because \(d\mapsto d\tanh(d\ell)\) is strictly
increasing; hence \(g_{d_2}/g_{d_1}\) is strictly decreasing in \(q\).  Thus
\(Q_{d_2}\) is smaller than \(Q_{d_1}\) in the likelihood-ratio order, and
therefore in the usual stochastic order \cite[Thm.~1.C.1]{ShakedShanthikumar2007}.  Since
both densities integrate to one and their ratio is strictly monotone, it
crosses \(1\) exactly once; hence \(g_{d_1}-g_{d_2}\) has a single sign change,
negative near \(0\) and positive near \(1/4\), so for every nondecreasing \(\phi\),
\[
 \E\phi(Q_{d_1})-\E\phi(Q_{d_2})
 =\int_0^{1/4}\bigl(\phi(q)-\phi(q^\ast)\bigr)
   \bigl(g_{d_1}(q)-g_{d_2}(q)\bigr)\,\mathrm d q\ge0,
\]
where \(q^\ast\) is the crossing point and the integrand is nonnegative
throughout.  Finally, symmetry of \(H\) gives \(H(p)=\widetilde H(p(1-p))\), with
\(\widetilde H\) nondecreasing on \([0,1/4]\) precisely when \(H\) is nondecreasing
on \([0,1/2]\); taking \(\phi=\widetilde H\) proves \eqref{eq:beta-order}.  If \(H\)
is strictly increasing and \(d_1<d_2\), the likelihood ratio is nonconstant, the
crossing is nondegenerate, and the inequality is strict.
\end{proof}

\section{The cubic residue kernel}\label{sec:kernel}

For symmetric nonnegative weights \(w_k=w_{m-k}\), the kernel of
\eqref{eq:G-weighted} is
\[
 G_{m,w}(p)=\sum_{k=1}^{m-1}w_k\binom{3m-2}{3k-1}
 p^{3k}(1-p)^{3(m-k)},
 \qquad G_{m,w}(p)=G_{m,w}(1-p).
\]
This section proves that \(G_{m,w}\) is monotone on \([0,1/2]\) whenever the
weights increase toward the center.

\begin{theorem}[Weighted cubic residue monotonicity]\label{thm:kernel}
Suppose
\begin{equation}\label{eq:w-monotone}
 0\le w_1\le w_2\le\cdots\le w_{\lfloor m/2\rfloor},
 \qquad w_k=w_{m-k}.
\end{equation}
Then \(G_{m,w}\) is nondecreasing on \([0,1/2]\).  If the weights are not all zero,
then \(G_{m,w}\) is strictly increasing on \([0,1/2]\).
\end{theorem}

\subsection{The constant-weight kernel}

Let \(w_k\equiv1\), and write \(G_m=G_{m,1}\).  Set \(n=3m-2\) and \(t=p/(1-p)\).
For \(0\le p\le1/2\) one has \(0\le t\le1\), and
\[
 G_m\!\left(\frac{t}{1+t}\right)
 =\frac{1}{(1+t)^{3m}}
 \sum_{k=1}^{m-1}\binom{n}{3k-1}t^{3k}.
\]
Differentiation gives
\begin{equation}\label{eq:G-J}
 \frac{\mathrm d}{\mathrm d t}G_m\!\left(\frac{t}{1+t}\right)
 =\frac{3J_m(t)}{(1+t)^{3m+1}},
\end{equation}
where
\begin{equation}\label{eq:J-k}
 J_m(t)=\sum_{k=1}^{m-1}\binom{n}{3k-1}t^{3k-1}
 \bigl(k-(m-k)t\bigr).
\end{equation}

\begin{lemma}\label{lem:bernstein}
For every \(m\ge2\),
\[
 J_m(t)>0\qquad(0<t<1).
\]
Consequently, \(G_m\) is strictly increasing on \([0,1/2]\).
\end{lemma}

\begin{proof}
Writing \(\nu=3k-1\) in \eqref{eq:J-k} gives
\[
 J_m(t)=\frac13
 \sum_{\substack{0\le\nu\le n\\\nu\equiv2\, (3)}}
 \binom{n}{\nu} t^\nu\bigl((\nu+1)-(n-\nu+1)t\bigr).
\]
Introduce the projective Bernstein transform
\begin{equation}\label{eq:P-def}
 \mathcal P_m(\xi)
 =(1+\xi)^{n+1}J_m\!\left(\frac{\xi}{1+\xi}\right).
\end{equation}
Then
\begin{equation}\label{eq:P-sum}
 \mathcal P_m(\xi)
 =\frac13
 \sum_{\nu\equiv2\, (3)}
 \binom{n}{\nu} \xi^\nu(1+\xi)^{n-\nu}
 \bigl((\nu+1)+(2\nu-n)\xi\bigr).
\end{equation}
For \(0\le j\le n+1\), coefficient extraction from \eqref{eq:P-sum} gives
\begin{align}
 [\xi^j]\mathcal P_m(\xi)
 &=\frac13
 \sum_{\substack{0\le\nu\le j\\\nu\equiv2\, (3)}}
 \binom{n}{\nu}
 \left[
 (\nu+1)\binom{n-\nu}{j-\nu}
 +(2\nu-n)\binom{n-\nu}{j-\nu-1}
 \right]\notag\\
 &=\frac{\binom{n+1}{j}}{3(n+1)}S_{n,j},
 \label{eq:P-coeff}
\end{align}
where
\begin{equation}\label{eq:S-def}
 S_{n,j}
 =\sum_{\substack{0\le\nu\le j\\\nu\equiv2\, (3)}}
 \binom{j}{\nu}(n-\nu+1)(2\nu+1-j).
\end{equation}
The second equality in \eqref{eq:P-coeff} follows from the elementary
identities
\[
 \binom{n}{\nu}\binom{n-\nu}{j-\nu}=\binom nj\binom{j}{\nu},
 \qquad
 \binom{n}{\nu}\binom{n-\nu}{j-\nu-1}
 =\binom{n+1}{j}\binom{j}{\nu}\frac{j-\nu}{n+1}.
\]
The transform is Bernstein-theoretic: writing
\(J_m(t)=\sum_j b_j\binom{n+1}{j}t^j(1-t)^{n+1-j}\) in the degree-\((n+1)\)
Bernstein basis, one has \([\xi^j]\mathcal P_m(\xi)=b_j\binom{n+1}{j}\), so
nonnegativity of the Bernstein coefficients \(b_j\) certifies \(J_m(t)\ge0\) on
\([0,1]\) \cite{PowersReznick2000}.  We claim that every \(S_{n,j}\) is nonnegative.

Put \(\delta=n+1-j\).  A third-root-of-unity filter evaluates the sum over the class
\(\nu\equiv2\pmod3\); its trigonometric terms are periodic modulo \(6\) in \(j\), so the
closed form splits into six residue classes,
\begin{equation}\label{eq:S-table}
3S_{n,j}=
\begin{cases}
 \delta(2^j-3j-1)-3j(j-1),&j\equiv0\pmod6,\\
 \delta(2^j-2)-3j(j-1),&j\equiv1\pmod6,\\
 \delta(2^j+3j-1),&j\equiv2\pmod6,\\
 \delta(2^j+3j+1)+3j(j-1),&j\equiv3\pmod6,\\
 \delta(2^j+2)+3j(j-1),&j\equiv4\pmod6,\\
 \delta(2^j-3j+1),&j\equiv5\pmod6.
\end{cases}
\end{equation}
In detail, for
\[
 R_a(j)=\sum_{\nu\equiv a\, (3)}\binom{j}{\nu}
\]
and \(\omega=e^{2\pi i/3}\),
\[
 R_a(j)=\frac13\sum_{\ell=0}^2\omega^{-a\ell}(1+\omega^\ell)^j
 =\frac13\left(2^j+2\cos\frac{(j-2a)\pi}{3}\right),
\]
where \(2\cos((j-2a)\pi/3)\) equals \(2,1,-1,-2,-1,1\) according as
\(j-2a\equiv0,1,2,3,4,5\pmod6\).  These values also follow with no appeal to
\(\mathbb C\): Pascal's rule gives \(R_a(j+1)=R_a(j)+R_{a-1}(j)\) (subscripts
modulo~\(3\)), and this recurrence with \(R_0(0)=1\), \(R_1(0)=R_2(0)=0\) fixes them
by induction.  The first two moments over the residue class are
\[
 \sum_{\nu\equiv2\, (3)}\nu\binom{j}{\nu}=jR_1(j-1),
 \qquad
 \sum_{\nu\equiv2\, (3)}\nu^2\binom{j}{\nu}
 =j(j-1)R_0(j-2)+jR_1(j-1).
\]
Expanding the summand in \eqref{eq:S-def} gives, for
\(j\ge2\),
\begin{align}
 S_{n,j}
 &=\delta\bigl(2jR_1(j-1)+(1-j)R_2(j)\bigr)\notag\\
 &\quad+j(j-1)\bigl(-2R_0(j-2)+3R_1(j-1)-R_2(j)\bigr),
 \label{eq:S-R}
\end{align}
with \(j=0,1\) following directly from \eqref{eq:S-def}.  Substituting the values
of \(R_2(j)\), \(R_1(j-1)\), and \(R_0(j-2)\) by residue \(j\bmod6\) yields
\eqref{eq:S-table}.

It remains only to inspect the six cases.  The cases
\(j\equiv2,3,4\pmod6\) are immediate.  If \(j\equiv5\pmod6\), then \(j\ge5\) and
\(2^j-3j+1>0\).
For the two remaining classes, note that \(n+1=3m-1\) is congruent to \(2\) or
\(5\) modulo \(6\); hence \(\delta\ge1\) when \(j\equiv1\pmod6\), and
\(\delta\ge2\) when \(j\equiv0\pmod6\) and \(j>0\).
If \(j\equiv1\pmod6\), the case \(j=1\) gives zero, while for
\(j\ge7\),
\begin{equation}\label{eq:exp-bound-1}
 2^j-2\ge3j(j-1).
\end{equation}
Finally, when \(j\equiv0\pmod6\) and \(j>0\), one has \(j\ge6\) and
\begin{equation}\label{eq:exp-bound-0}
 2(2^j-3j-1)\ge3j(j-1).
\end{equation}
The inequalities \eqref{eq:exp-bound-1} and \eqref{eq:exp-bound-0} hold with
equality at \(j=7\) and \(j=6\), respectively.  Advancing \(j\) by one increases the
left-minus-right difference by \(2^j-6j\) in \eqref{eq:exp-bound-1} and by
\(2^{j+1}-6j-6\) in \eqref{eq:exp-bound-0}, both positive for \(j\ge6\); the
inequalities therefore persist for all larger \(j\).  Likewise \(2^j-3j+1>0\) at
\(j=5\), with increment \(2^j-3>0\) thereafter.  Hence \(S_{n,j}\ge0\) in every case.

By \eqref{eq:P-coeff}, all coefficients of \(\mathcal P_m\) are nonnegative.
The coefficient at \(\xi^2\) is positive, since \(\delta=n-1>0\) and
\(3S_{n,2}=9\delta\).  Therefore \(\mathcal P_m(\xi)>0\) for \(\xi>0\).  The substitution
\(t=\xi/(1+\xi)\) in \eqref{eq:P-def} proves \(J_m(t)>0\) for \(0<t<1\), and
\eqref{eq:G-J} completes the proof.
\end{proof}

\subsection{Insertion of monotone weights}

Arbitrary weights satisfying \eqref{eq:w-monotone} preserve nonnegativity of the
derivative.

For \(0<t<1\), differentiation of \eqref{eq:G-weighted} gives
\begin{equation}\label{eq:J-weighted}
 \frac{\mathrm d}{\mathrm d t}G_{m,w}\!\left(\frac{t}{1+t}\right)
 =\frac{3J_{m,w}(t)}{(1+t)^{3m+1}},
\end{equation}
where
\begin{equation}\label{eq:Jw-def}
 J_{m,w}(t)
 =\sum_{k=1}^{m-1}w_k\binom{3m-2}{3k-1}t^{3k-1}
 \bigl(k-(m-k)t\bigr).
\end{equation}
Pair the terms indexed by \(k\) and \(m-k\).  For \(1\le k<m/2\), let
\begin{align}
 B_{m,k}(t)
 &=\binom{3m-2}{3k-1}
 \Bigl[t^{3k-1}\bigl(k-(m-k)t\bigr)\notag\\
 &\hspace{36mm}
 +t^{3(m-k)-1}\bigl((m-k)-kt\bigr)\Bigr].
 \label{eq:B-def}
\end{align}
If \(m\) is even, the central block is
\begin{equation}\label{eq:B-center}
 B_{m,m/2}(t)
 =\binom{3m-2}{3m/2-1}\frac m2\,t^{3m/2-1}(1-t)>0.
\end{equation}
Then
\[
 J_{m,w}(t)
 =\sum_{k=1}^{\lfloor m/2\rfloor}w_kB_{m,k}(t),
\]
with the convention \eqref{eq:B-center} at the center.

\begin{lemma}\label{lem:block-sign}
For each fixed \(m\ge2\) and \(0<t<1\), the sequence
\[
 B_{m,1}(t),B_{m,2}(t),\dots,B_{m,\lfloor m/2\rfloor}(t)
\]
has no more than one sign change, from negative to positive, and its final term
\(B_{m,\lfloor m/2\rfloor}(t)\) is positive.
\end{lemma}

\begin{proof}
For \(k<m/2\), the sign of \(B_{m,k}(t)\) is the sign of
\[
 H_{m,k}(t)
 =k-(m-k)t+t^{3(m-2k)}\bigl((m-k)-kt\bigr).
\]
Put \(d=m-2k>0\) and \(t=e^{-z}\), \(z>0\).  Direct rearrangement gives
\begin{align}
 H_{m,k}(t)
 &=\frac12(1+t)(1+t^{3d})
 \left[
 m\frac{1-t}{1+t}
 -d\frac{1-t^{3d}}{1+t^{3d}}
 \right]\notag\\
 &=\frac12(1+t)(1+t^{3d})
 \left[m\tanh\frac{z}{2}-d\tanh\frac{3dz}{2}\right].
 \label{eq:H-hyperbolic}
\end{align}
The function \(d\mapsto d\tanh(3dz/2)\) is strictly increasing for \(d>0\).
Therefore, for \(2\le k<m/2\), if
\(B_{m,k}(t)\le0\), then the same bracket in \eqref{eq:H-hyperbolic} is strictly
negative after replacing \(d\) by \(d+2\), which corresponds to replacing \(k\) by
\(k-1\).  Thus
\[
 B_{m,k}(t)\le0\quad\Longrightarrow\quad B_{m,k-1}(t)<0
 \qquad(2\le k<m/2).
\]
The final block is always positive.  For even \(m\) this is \eqref{eq:B-center}.
For odd \(m\), the final block has \(d=1\), and its bracket in
\eqref{eq:H-hyperbolic} is \(m\tanh(z/2)-\tanh(3z/2)>0\),
because \(m\ge3\) and \(\tanh(3y)<3\tanh y\) for \(y>0\).  Hence the blocks have the
asserted sign pattern.
\end{proof}

\begin{proof}[Proof of \Cref{thm:kernel}]
By \Cref{lem:bernstein},
\[
 \sum_{k=1}^{\lfloor m/2\rfloor}B_{m,k}(t)=J_m(t)>0
 \qquad(0<t<1).
\]
The block sequence has one sign change by \Cref{lem:block-sign}.  Applying
\Cref{lem:weighting} with the nondecreasing weights in
\eqref{eq:w-monotone} gives
\[
 J_{m,w}(t)=\sum_{k=1}^{\lfloor m/2\rfloor}w_kB_{m,k}(t)\ge0.
\]
Equation \eqref{eq:J-weighted} proves that \(G_{m,w}\) is nondecreasing on
\([0,1/2]\).

Suppose now that the weights are not all zero, and put \(L=\lfloor m/2\rfloor\).
Then \(w_L>0\), and \(B_{m,L}(t)>0\) by \Cref{lem:block-sign}.  Since
\(\sum_{k=1}^LB_{m,k}(t)=J_m(t)>0\), the strict part of \Cref{lem:weighting} gives
\(J_{m,w}(t)>0\) for \(0<t<1\), and \(G_{m,w}\) is strictly increasing on
\([0,1/2]\).
\end{proof}

\section{Proof of the theorem}\label{sec:proof}

The representation \eqref{eq:C-beta} and the two monotonicity facts of
\Cref{sec:monotonicity-lemmas,sec:kernel} combine to show that the coefficient
convolution is Schur-concave in its two parameters
\cite[Def.~3.A.1]{MarshallOlkinArnold2011}, which gives \Cref{thm:main}.

\begin{proposition}
\label{prop:C-schur}
For each \(m\ge2\) and each log-concave sequence \((f_n)\) as in
\Cref{thm:main}, the symmetric function \((u,v)\mapsto C_{m,f}(u,v)\) is
Schur-concave on \((0,\infty)^2\).  Equivalently, at fixed \(s=u+v\),
\[
 d\longmapsto
 C_{m,f}\left(\frac s2+d,\frac s2-d\right)
\]
is nonincreasing for \(0\le d<s/2\).  If at least one product \(f_kf_{m-k}\) is
positive, then this function is strictly decreasing.
\end{proposition}

\begin{proof}
The prefactor in \eqref{eq:C-beta} depends only on \(s=u+v\), and \(G_{m,w}\) is
symmetric by \eqref{eq:G-weighted}.  By \Cref{thm:kernel} it is nondecreasing on
\([0,1/2]\), and strictly increasing when at least one weight \(w_k=f_kf_{m-k}\)
is positive; \Cref{lem:beta-order} and its strict part give the two conclusions.
\end{proof}

\begin{proof}[Proof of \Cref{thm:main}]
The coefficients of degrees \(0\) and \(1\) in \eqref{eq:Turan-def} vanish.  For
\(m\ge2\), direct multiplication gives
\begin{align}
 [x^m]\Tur_f(\mu;\alpha,\beta;x)
 &=C_{m,f}(\mu+\alpha,\mu+\beta)\notag\\
 &\quad-C_{m,f}(\mu,\mu+\alpha+\beta).
 \label{eq:delta-C}
\end{align}
Assume first that \(\mu>0\) and \(\alpha,\beta>0\).  The two parameter pairs in
\eqref{eq:delta-C} have the same sum \(s=2\mu+\alpha+\beta\) and imbalances
\[
 d_1=\frac{|\alpha-\beta|}{2}\le\frac{\alpha+\beta}{2}=d_2.
\]
\Cref{prop:C-schur} therefore gives
\[
 C_{m,f}(\mu+\alpha,\mu+\beta)
 \ge C_{m,f}(\mu,\mu+\alpha+\beta).
\]
This proves nonnegativity of every coefficient.  If \(\alpha=0\) or \(\beta=0\),
the Tur\'anian is identically zero.  The boundary case \(\mu=0\) follows from
\(F_f(0;x)=0\) and the nonnegative coefficients of
\(F_f(\alpha;x)F_f(\beta;x)\).
\end{proof}

\begin{corollary}\label{cor:strict}
Under the hypotheses of \Cref{thm:main}, let \(\mu\ge0\), \(\alpha,\beta>0\), and
\(m\ge2\).  Then
\[
 [x^m]\Tur_f(\mu;\alpha,\beta;x)>0
\]
if and only if there exists \(1\le k\le m-1\) with \(f_kf_{m-k}>0\).
\end{corollary}

\begin{proof}
If every product \(f_kf_{m-k}\) vanishes, then both terms of \eqref{eq:delta-C}
vanish.  Conversely, suppose one such product is positive.  If \(\mu>0\), the two
imbalances in the proof of \Cref{thm:main} satisfy \(d_1<d_2\), so the strict
part of \Cref{prop:C-schur} gives strict inequality.  If \(\mu=0\),
then \(F_f(0;x)=0\), so the coefficient equals \(C_{m,f}(\alpha,\beta)\), a sum of
nonnegative terms with at least one strictly positive.
\end{proof}

\begin{corollary}\label{cor:ordinary}
Suppose \((f_n)\) is not identically zero, and let \(R\) be the radius of
convergence of \(\sum_{n\ge1}f_nx^n\).  For every \(0<x<R\), the function
\(\mu\mapsto F_f(\mu;x)\) is finite, continuous, and positive on \((0,\infty)\), and
\(\mu\mapsto\log F_f(\mu;x)\) is concave there.
\end{corollary}

\begin{proof}
Log-concavity with no internal zeros makes \(f_{n+1}/f_n\) nonincreasing on the
support of \((f_n)\), so \(R>0\).  For each fixed \(\mu>0\), Stirling's formula gives
\[
 \frac{(\mu)_{3n}}{(3n-1)!}
 =\frac{\Gamma(3n+\mu)}{\Gamma(\mu)\Gamma(3n)}
 \sim\frac{(3n)^\mu}{\Gamma(\mu)}
 \qquad(n\to\infty).
\]
Hence \(F_f(\mu;\cdot)\) has the same radius of convergence \(R\) as
\(\sum_{n\ge1}f_nx^n\).  On every compact set of positive \(\mu\) and for
\(0<x<R\), the series converges locally uniformly in \(\mu\), since its extra
factor grows at most polynomially in \(n\); thus \(\mu\mapsto F_f(\mu;x)\) is
continuous.  Since \((f_n)\) is not identically zero, its nonnegative
coefficients give \(F_f(\mu;x)>0\) for \(\mu,x>0\).  By \Cref{thm:main},
\[
 F_f(\mu+\alpha;x)F_f(\mu+\beta;x)
 \ge F_f(\mu;x)F_f(\mu+\alpha+\beta;x)
\]
for all nonnegative shifts.  Taking \(\alpha=\beta=h>0\) gives
\(F_f(\mu+h;x)^2\ge F_f(\mu;x)F_f(\mu+2h;x)\),
so \(\mu\mapsto\log F_f(\mu;x)\) is midpoint concave, and by continuity concave,
on \((0,\infty)\).
\end{proof}

\begin{remark}\label{rem:internal-zeros}
The interval-support condition cannot be dropped.  Put \(f_1=f_4=1\) and
\(f_n=0\) otherwise.  The two positive indices differ by \(3\), so
\(f_{n-1}f_{n+1}=0\) at every index where the three terms are defined and
\(f_n^2\ge f_{n-1}f_{n+1}\) holds throughout, strictly at \(n=4\).  The positive
support \(\{1,4\}\) is not an interval.  A single zero between two positive terms
violates the inequality at that zero, so a gap of two is the shortest the
inequality admits.  At \(m=5\) the weights \(w_k=f_kf_{5-k}\) are
\(w_1=w_4=1\) and \(w_2=w_3=0\), decreasing toward the center rather than
increasing, and
\[
 [x^5]\Tur_f(2;1,1;x)=-9360.
\]
\end{remark}

\section{Concluding remarks}

The proof of \Cref{thm:main} hinges on the special feature of the residue class
modulo \(3\).  After the beta-binomial reduction, the entire conjecture is
controlled by the monotonicity of
\[
 p(1-p)\Pp\{\operatorname{Bin}(3m-2,p)\equiv2\pmod3\}.
\]
The projective transform \eqref{eq:P-def} converts its derivative into a
polynomial whose Bernstein coefficients are nonnegative.  This positivity is
not termwise obvious in the monomial basis; it appears only after the
root-of-unity filter and its mod-\(6\) case split.  Log-concavity enters
the proof of \Cref{thm:main} exactly once: it
makes the symmetric weights \(w_k=f_kf_{m-k}\) nondecreasing as \(k\) approaches
\(m/2\).

For a general integer \(r\), the same beta-binomial reduction produces, in the
constant-weight case,
\[
 p(1-p)\Pp\{\operatorname{Bin}(rm-2,p)\equiv r-1\pmod r\},
\]
and in general its analogue weighted by the products \(f_kf_{m-k}\).
The cubic proof suggests testing whether a suitable projective degree lift has
a positive coefficient array.  The known failure at \(r=4\) shows that such
positivity cannot persist without an additional restriction; locating the
first coefficient obstruction in the lifted array may explain that failure.

\section*{Acknowledgements}
During the preparation of this work, the author used OpenAI's ChatGPT (GPT-5.6 Sol) and Anthropic's Claude (Opus 5, Opus 4.8) to assist with manuscript drafting and the development of auxiliary verification code.

\section*{Statements and declarations}

\noindent\textbf{Funding.} This research received no external funding.

\smallskip
\noindent\textbf{Competing interests.} The author has no relevant interests to disclose.

\smallskip
\noindent\textbf{Data availability.} No datasets were generated or analysed.

\smallskip
\noindent\textbf{Computational verification.} The proofs do not depend on numerical computation.  Scripts used for verification are available at \url{https://github.com/johnnyshields/papers}.

\begin{thebibliography}{99}

\bibitem{Baricz2008}
\'A. Baricz,
\emph{Tur\'an type inequalities for hypergeometric functions},
Proc. Amer. Math. Soc. \textbf{136} (2008), no.~9, 3223--3229,
DOI: \href{https://doi.org/10.1090/S0002-9939-08-09353-2}{10.1090/S0002-9939-08-09353-2}.

\bibitem{BarnardGordyRichards2009}
R.~W. Barnard, M.~B. Gordy, and K.~C. Richards,
\emph{A note on Tur\'an type and mean inequalities for the Kummer function},
J. Math. Anal. Appl. \textbf{349} (2009), no.~1, 259--263,
DOI: \href{https://doi.org/10.1016/j.jmaa.2008.08.024}{10.1016/j.jmaa.2008.08.024}.

\bibitem{BerenhautBergen2011}
K.~S. Berenhaut and L.~D. Bergen,
\emph{Stochastic orderings, folded beta distributions and fairness in coin flips},
Statist. Probab. Lett. \textbf{81} (2011), no.~6, 632--638,
DOI: \href{https://doi.org/10.1016/j.spl.2011.01.002}{10.1016/j.spl.2011.01.002}.

\bibitem{KalmykovKarp2013Reciprocal}
S.~I. Kalmykov and D.~B. Karp,
\emph{Log-concavity for series in reciprocal gamma functions and applications},
Integral Transforms Spec. Funct. \textbf{24} (2013), no.~11, 859--872,
DOI: \href{https://doi.org/10.1080/10652469.2013.764874}{10.1080/10652469.2013.764874}.

\bibitem{KalmykovKarp2013Ratios}
S.~I. Kalmykov and D.~B. Karp,
\emph{Log-convexity and log-concavity for series in gamma ratios and applications},
J. Math. Anal. Appl. \textbf{406} (2013), no.~2, 400--418,
DOI: \href{https://doi.org/10.1016/j.jmaa.2013.04.061}{10.1016/j.jmaa.2013.04.061}.

\bibitem{KarpSitnik2010}
D.~Karp and S.~M. Sitnik,
\emph{Log-convexity and log-concavity of hypergeometric-like functions},
J. Math. Anal. Appl. \textbf{364} (2010), no.~2, 384--394,
DOI: \href{https://doi.org/10.1016/j.jmaa.2009.10.057}{10.1016/j.jmaa.2009.10.057}.

\bibitem{KarpZhang2024}
D.~Karp and Y.~Zhang,
\emph{Log-concavity and log-convexity of series containing multiple Pochhammer symbols},
Fract. Calc. Appl. Anal. \textbf{27} (2024), no.~1, 458--486,
DOI: \href{https://doi.org/10.1007/s13540-023-00238-0}{10.1007/s13540-023-00238-0}.

\bibitem{MarshallOlkinArnold2011}
A.~W. Marshall, I.~Olkin, and B.~C. Arnold,
\emph{Inequalities: Theory of Majorization and Its Applications},
2nd ed., Springer Series in Statistics, Springer, New York, 2011,
DOI: \href{https://doi.org/10.1007/978-0-387-68276-1}{10.1007/978-0-387-68276-1}.

\bibitem{NinhPham2025}
A.~Ninh and M.~Pham,
\emph{Logconcavity, twice-logconcavity and Tur\'an-type inequalities},
Ann. Oper. Res. \textbf{354} (2025), no.~3, 1271--1283,
DOI: \href{https://doi.org/10.1007/s10479-018-2923-y}{10.1007/s10479-018-2923-y}.

\bibitem{PowersReznick2000}
V.~Powers and B.~Reznick,
\emph{Polynomials that are positive on an interval},
Trans. Amer. Math. Soc. \textbf{352} (2000), no.~10, 4677--4692,
DOI: \href{https://doi.org/10.1090/S0002-9947-00-02595-2}{10.1090/S0002-9947-00-02595-2}.

\bibitem{ShakedShanthikumar2007}
M.~Shaked and J.~G. Shanthikumar,
\emph{Stochastic Orders},
Springer Series in Statistics, Springer, New York, 2007,
DOI: \href{https://doi.org/10.1007/978-0-387-34675-5}{10.1007/978-0-387-34675-5}.

\end{thebibliography}

\end{document}
